% =============================================================================
% APPENDIX C: CORE-WHAT: Utility Dimensions (FEPSDE)
% =============================================================================
% Module: BCM2_04_INU (Utility Dimensions)
% Category: CORE
% Version: 1.0 (January 2026)
% Status: Final
% Language: English
%
% PURPOSE: This appendix answers the fundamental question "What constitutes utility?"
% It formalizes the FEPSDE taxonomy: Financial, Experiential, Psychological, Social,
% Developmental, and Existential utility dimensions.
%
% 10C POSITION: This is CORE 2 of 10 - the "WHAT" question.
%
% =============================================================================

% -----------------------------------------------------------------------------
% HEADER BLOCK
% -----------------------------------------------------------------------------

\begin{tcolorbox}[colback=gray!5!white, colframe=gray!75!black,
    title=Appendix Header: C CORE-WHAT]

\begin{tabular}{@{}ll@{}}
\textbf{Code:} & C \\
\textbf{Category:} & CORE (Fundamental Theory) \\
\textbf{Full Name:} & CORE-WHAT: Utility Dimensions (FEPSDE) \\
\textbf{Module:} & BCM2\_04\_INU (Utility Dimensions) \\
\textbf{Version:} & 1.0 (January 2026) \\
\textbf{Status:} & Final \\
\textbf{Language:} & English \\
\end{tabular}

\vspace{0.5em}
\textbf{Purpose:} Formalizes the \textbf{FEPSDE Taxonomy}, answering the
second fundamental question of the EBF: \textit{``What constitutes utility?''}
Establishes six irreducible utility dimensions: Financial (F), Experiential (E),
Psychological (P), Social (S), Developmental (D), and Existential (X).

\textbf{10C Position:} This is CORE 2 (WHAT) --- builds on WHO (AAA),
provides content for HOW (B) interactions.

\end{tcolorbox}

% =============================================================================
% CROSS-REFERENCE STRUCTURE
% =============================================================================

\begin{tcolorbox}[colback=purple!5!white, colframe=purple!75!black,
    title=Chapter Linkage]

\textbf{Primary Chapter:} Chapter 6: Utility Dimensions

\vspace{0.3em}
\textbf{The Fundamental Question:}
\begin{quote}
\textit{``What constitutes utility? What are the irreducible dimensions of human welfare?''}
\end{quote}

\textbf{Section-Level Mapping:}
\begin{itemize}[nosep]
\item Section 6.1: Beyond Money $\leftrightarrow$ Section~\ref{sec:what-fundamental}
\item Section 6.2: The FEPSDE Taxonomy $\leftrightarrow$ Section~\ref{sec:what-fepsde}
\item Section 6.3: Measurement $\leftrightarrow$ Section~\ref{sec:what-measurement}
\item Section 6.4: Integration $\leftrightarrow$ Section~\ref{sec:what-integration}
\end{itemize}

\textbf{Secondary Chapters:}
\begin{itemize}[nosep]
\item Chapter 3: Utility Foundations --- Mathematical basis
\item Chapter 5: Welfare Hierarchy --- Whose utility dimensions
\item Chapter 7: Complementarity --- Cross-dimension interactions
\end{itemize}

\textbf{Relationship to Other COREs:}
\begin{itemize}[nosep]
\item Appendix UNMAPPED_AAA (CORE-WHO): Provides bearers whose utility dimensions WHAT specifies
\item Appendix UNMAPPED_B (CORE-HOW): Uses WHAT dimensions for complementarity matrix
\item Appendix UNMAPPED_V (CORE-WHEN): Context affects dimension weights
\end{itemize}

\end{tcolorbox}

% -----------------------------------------------------------------------------
% APPENDIX SCOPE BOX
% -----------------------------------------------------------------------------

\begin{tcolorbox}[
    colback=orange!5!white,
    colframe=orange!75!black,
    title={\textbf{Appendix Scope: Ziel / In-Scope / Out-of-Scope / Constraints}},
    fonttitle=\bfseries
]

\textbf{Ziel (Objective):}
\begin{itemize}[nosep]
    \item Define the complete taxonomy of utility dimensions (FEPSDE)
\end{itemize}

\textbf{In-Scope:}
\begin{itemize}[nosep]
    \item Six dimension definitions (F, E, P, S, D, X)
    \item Axioms DIM-1 through DIM-12
    \item Measurement approaches per dimension
    \item Sub-dimensions and hierarchical structure
\end{itemize}

\textbf{Out-of-Scope:}
\begin{itemize}[nosep]
    \item Whose utility $\rightarrow$ Appendix UNMAPPED_AAA (CORE-WHO)
    \item Cross-dimension interactions $\rightarrow$ Appendix UNMAPPED_B (CORE-HOW)
    \item Context-dependent weighting $\rightarrow$ Appendix UNMAPPED_V (CORE-WHEN)
    \item Empirical parameter values $\rightarrow$ Appendix UNMAPPED_BBB (CORE-WHERE)
\end{itemize}

\textbf{Constraints:}
\begin{itemize}[nosep]
    \item Dimensions must be mutually exclusive at top level
    \item Each dimension must be empirically measurable
    \item Taxonomy must be culturally general
\end{itemize}

\textbf{Lieferobjekte (Deliverables):}
\begin{itemize}[nosep]
    \item Complete FEPSDE Taxonomy (Table~\ref{tab:what-fepsde})
    \item Axiom System DIM-1 through DIM-12
    \item Measurement Matrix per dimension
    \item Utility aggregation formula (Equation~\ref{eq:what-utility})
\end{itemize}

\end{tcolorbox}

% =============================================================================
% CHAPTER
% =============================================================================

\chapter{CORE-WHAT: Utility Dimensions (FEPSDE)}
\label{app:core-what}

\begin{abstract}
This appendix answers the second fundamental question of the Evidence-Based Framework:
\textit{``What constitutes utility?''} Standard economics reduces utility to a single
dimension (typically money or consumption). We establish that human welfare comprises
six irreducible dimensions: Financial (F), Experiential (E), Psychological (P),
Social (S), Developmental (D), and Existential (X) --- the FEPSDE taxonomy.
The axiom system DIM-1 through DIM-12 formalizes these dimensions, their properties,
and measurement approaches.
\end{abstract}

% -----------------------------------------------------------------------------
% QUICK REFERENCE BOX
% -----------------------------------------------------------------------------

\begin{tcolorbox}[colback=blue!5!white, colframe=blue!75!black,
    title=Quick Reference: CORE-WHAT]

\textbf{Core Question:} What constitutes utility? What are the dimensions of welfare?

\textbf{Key Formula:}
\begin{equation}
U = \sum_{d \in \{F,E,P,S,D,X\}} \alpha_d \cdot u_d + \sum_{d < d'} \gamma_{dd'} \cdot u_d \cdot u_{d'}
\label{eq:what-utility}
\end{equation}

\textbf{The FEPSDE Dimensions:}
\begin{itemize}[nosep]
\item \textbf{F}inancial: Material resources, wealth, income
\item \textbf{E}xperiential: Hedonic pleasure, positive emotions
\item \textbf{P}sychological: Autonomy, competence, mental health
\item \textbf{S}ocial: Belonging, relationships, status
\item \textbf{D}evelopmental: Growth, learning, self-improvement
\item e\textbf{X}istential: Meaning, purpose, transcendence
\end{itemize}

\textbf{Cross-References:}
AAA (CORE-WHO), B (CORE-HOW), V (CORE-WHEN), AU (CORE-AWARE), CP (DOMAIN-CPS: U\_D Kompetenz)

\end{tcolorbox}


%%%%%%%%%%%%%%%%%%%%%%%%%%%%%%%%%%%%%%%%%%%%%%%%%%%%%%%%%%%%%%%%%%%%%%%%%%%%%%%
\section{The Fundamental Question}
\label{sec:what-fundamental}
%%%%%%%%%%%%%%%%%%%%%%%%%%%%%%%%%%%%%%%%%%%%%%%%%%%%%%%%%%%%%%%%%%%%%%%%%%%%%%%

\subsection{Why "What" Matters}

Economics traditionally assumes utility is unidimensional --- more consumption
is better. But this misses essential aspects of human welfare:

\begin{itemize}
\item A wealthy person may be miserable (low E, P)
\item A poor community may have high social capital (high S)
\item Material abundance without meaning feels empty (low X)
\item Growth and learning matter beyond hedonic pleasure (D $\neq$ E)
\end{itemize}

\subsection{The Core Problem}

Unidimensional utility models create blind spots:

\begin{enumerate}
\item \textbf{Policy Myopia}: Focus on GDP ignores non-material welfare
\item \textbf{Measurement Gaps}: What cannot be measured in money is ignored
\item \textbf{Intervention Failures}: Policies that boost one dimension may
      harm others unintentionally
\end{enumerate}

\begin{tcolorbox}[colback=red!5!white, colframe=red!75!black,
    title=The Fundamental Question]
\textbf{What constitutes utility? What are the irreducible dimensions of human welfare that any complete framework must include?}
\end{tcolorbox}


%%%%%%%%%%%%%%%%%%%%%%%%%%%%%%%%%%%%%%%%%%%%%%%%%%%%%%%%%%%%%%%%%%%%%%%%%%%%%%%
\section{Core Theory: The FEPSDE Taxonomy}
\label{sec:what-fepsde}
%%%%%%%%%%%%%%%%%%%%%%%%%%%%%%%%%%%%%%%%%%%%%%%%%%%%%%%%%%%%%%%%%%%%%%%%%%%%%%%

\subsection{The Six Dimensions}

\begin{table}[htbp]
\centering
\caption{The FEPSDE Utility Taxonomy}
\label{tab:what-fepsde}
\renewcommand{\arraystretch}{1.4}
\begin{tabular}{@{}clp{5cm}l@{}}
\toprule
\textbf{Code} & \textbf{Dimension} & \textbf{Definition} & \textbf{Theoretical Basis} \\
\midrule
F & Financial & Material resources, wealth, income, security & Classical economics \\
E & Experiential & Hedonic pleasure, positive emotions, absence of pain & Kahneman's experienced utility \\
P & Psychological & Autonomy, competence, relatedness, mental health & Self-Determination Theory \\
S & Social & Belonging, relationships, status, recognition & Social identity theory \\
D & Developmental & Growth, learning, skill acquisition, self-improvement & Maslow, Ryff \\
X & Existential & Meaning, purpose, values, transcendence & Frankl, eudaimonic well-being \\
\bottomrule
\end{tabular}
\end{table}

\subsection{Dimension Definitions}

\subsubsection{F: Financial Utility}

\begin{definition}[Financial Utility]
\textbf{Financial utility} $u_F$ captures welfare derived from material resources:
wealth, income, financial security, consumption capacity, and asset accumulation.
\end{definition}

\textbf{Sub-dimensions:}
\begin{itemize}[nosep]
\item $u_{F,\text{income}}$: Regular income flows
\item $u_{F,\text{wealth}}$: Asset stocks
\item $u_{F,\text{security}}$: Insurance against loss
\item $u_{F,\text{liquidity}}$: Access to funds when needed
\end{itemize}

\textbf{Measurement:} Objective (income, assets) and subjective (financial satisfaction).

\subsubsection{E: Experiential Utility}

\begin{definition}[Experiential Utility]
\textbf{Experiential utility} $u_E$ captures moment-to-moment hedonic experience:
pleasure, positive emotions, absence of pain and suffering.
\end{definition}

\textbf{Sub-dimensions:}
\begin{itemize}[nosep]
\item $u_{E,\text{pleasure}}$: Positive affect
\item $u_{E,\text{pain}}$: Absence of negative affect (inverted)
\item $u_{E,\text{flow}}$: Optimal engagement states
\item $u_{E,\text{sensory}}$: Aesthetic and sensory enjoyment
\end{itemize}

\textbf{Measurement:} Experience Sampling Method (ESM), Day Reconstruction Method.

\subsubsection{P: Psychological Utility}

\begin{definition}[Psychological Utility]
\textbf{Psychological utility} $u_P$ captures welfare from psychological need
satisfaction: autonomy, competence, relatedness, and mental health.
\end{definition}

\textbf{Sub-dimensions (SDT-aligned):}
\begin{itemize}[nosep]
\item $u_{P,\text{autonomy}}$: Self-determination, volition
\item $u_{P,\text{competence}}$: Mastery, efficacy
\item $u_{P,\text{relatedness}}$: Connection (overlaps with S)
\item $u_{P,\text{health}}$: Absence of mental illness
\end{itemize}

\textbf{Measurement:} Basic Psychological Need Satisfaction Scale, mental health inventories.

\subsubsection{S: Social Utility}

\begin{definition}[Social Utility]
\textbf{Social utility} $u_S$ captures welfare from social relationships:
belonging, love, friendship, status, recognition, and social identity.
\end{definition}

\textbf{Sub-dimensions:}
\begin{itemize}[nosep]
\item $u_{S,\text{belonging}}$: Group membership
\item $u_{S,\text{intimacy}}$: Close relationships
\item $u_{S,\text{status}}$: Social rank, prestige
\item $u_{S,\text{recognition}}$: Acknowledgment by others
\end{itemize}

\textbf{Measurement:} Social network analysis, social support scales, status measures.

\subsubsection{D: Developmental Utility}

\begin{definition}[Developmental Utility]
\textbf{Developmental utility} $u_D$ captures welfare from growth and learning:
skill acquisition, personal development, self-improvement, and becoming.
\end{definition}

\textbf{Sub-dimensions:}
\begin{itemize}[nosep]
\item $u_{D,\text{learning}}$: Skill/knowledge acquisition
\item $u_{D,\text{growth}}$: Personal development trajectory
\item $u_{D,\text{potential}}$: Future capability expansion
\item $u_{D,\text{mastery}}$: Achievement of expertise
\end{itemize}

\textbf{Measurement:} Personal growth inventories, skill assessments, trajectory analysis.

\subsubsection{X: Existential Utility}

\begin{definition}[Existential Utility]
\textbf{Existential utility} $u_X$ captures welfare from meaning and purpose:
life meaning, values alignment, transcendence, and spiritual fulfillment.
\end{definition}

\textbf{Sub-dimensions:}
\begin{itemize}[nosep]
\item $u_{X,\text{meaning}}$: Sense of life significance
\item $u_{X,\text{purpose}}$: Goal-directedness
\item $u_{X,\text{values}}$: Living according to values
\item $u_{X,\text{transcendence}}$: Connection to something larger
\end{itemize}

\textbf{Measurement:} Meaning in Life Questionnaire, Purpose scales, values congruence.


%%%%%%%%%%%%%%%%%%%%%%%%%%%%%%%%%%%%%%%%%%%%%%%%%%%%%%%%%%%%%%%%%%%%%%%%%%%%%%%
\section{Axioms and Formal Definitions}
\label{sec:what-axioms}
%%%%%%%%%%%%%%%%%%%%%%%%%%%%%%%%%%%%%%%%%%%%%%%%%%%%%%%%%%%%%%%%%%%%%%%%%%%%%%%

% DIM-1
\begin{tcolorbox}[colback=gray!5!white, colframe=gray!75!black,
    title=Axiom DIM-1: Multi-Dimensionality]
\textbf{Statement:}
\begin{equation}
U = f(u_F, u_E, u_P, u_S, u_D, u_X) \quad \text{where } \frac{\partial f}{\partial u_d} > 0 \; \forall d
\end{equation}

\textbf{Interpretation:} Utility is multi-dimensional. Each dimension contributes
positively to total utility, ceteris paribus.

\textbf{Implication:} No single dimension captures all of welfare.
\end{tcolorbox}

% DIM-2
\begin{tcolorbox}[colback=gray!5!white, colframe=gray!75!black,
    title=Axiom DIM-2: Irreducibility]
\textbf{Statement:}
\begin{equation}
\nexists \; d, d': u_d = g(u_{d'}) \quad \text{for all entities}
\end{equation}

\textbf{Interpretation:} No dimension is fully reducible to another.
Financial wealth cannot substitute for meaning; pleasure cannot replace growth.

\textbf{Implication:} All six dimensions are necessary; none is redundant.
\end{tcolorbox}

% DIM-3
\begin{tcolorbox}[colback=gray!5!white, colframe=gray!75!black,
    title=Axiom DIM-3: Boundedness]
\textbf{Statement:}
\begin{equation}
\forall d: u_d \in [0, 1] \quad \text{(normalized scale)}
\end{equation}

\textbf{Interpretation:} Each dimension is bounded and normalized to [0, 1]
for cross-dimension comparability.

\textbf{Implication:} Enables meaningful aggregation across dimensions.
\end{tcolorbox}

% DIM-4
\begin{tcolorbox}[colback=gray!5!white, colframe=gray!75!black,
    title=Axiom DIM-4: Diminishing Marginal Utility]
\textbf{Statement:}
\begin{equation}
\frac{\partial^2 U}{\partial u_d^2} < 0 \quad \text{(within dimension)}
\end{equation}

\textbf{Interpretation:} Each dimension exhibits diminishing marginal utility
within itself. The 10th dollar matters less than the first.

\textbf{Exception:} Between dimensions, non-concavity possible via $\gamma_{dd'} > 0$.
\end{tcolorbox}

% DIM-5
\begin{tcolorbox}[colback=gray!5!white, colframe=gray!75!black,
    title=Axiom DIM-5: Universal Presence]
\textbf{Statement:}
\begin{equation}
\forall \text{ cultures, } \forall \text{ levels } L: \text{FEPSDE applies}
\end{equation}

\textbf{Interpretation:} The six dimensions are universal across cultures and
welfare levels. The \textit{weights} may vary, but the \textit{dimensions} are constant.

\textbf{Empirical Basis:} Cross-cultural well-being research; universal psychological needs.
\end{tcolorbox}

% DIM-6
\begin{tcolorbox}[colback=gray!5!white, colframe=gray!75!black,
    title=Axiom DIM-6: Weight Heterogeneity]
\textbf{Statement:}
\begin{equation}
\alpha_d = \alpha_d(i, L, \Psi) \quad \text{varies by individual, level, context}
\end{equation}

\textbf{Interpretation:} The importance weight of each dimension varies across
individuals, welfare levels, and contexts.

\textbf{Cross-Reference:} Weight determination via CORE-WHEN (V) and CORE-WHERE (BBB).
\end{tcolorbox}

% DIM-7
\begin{tcolorbox}[colback=gray!5!white, colframe=gray!75!black,
    title=Axiom DIM-7: Measurability]
\textbf{Statement:}
\begin{equation}
\forall d: \exists \text{ valid, reliable measure } M_d
\end{equation}

\textbf{Interpretation:} Each dimension must be empirically measurable.
Dimensions without valid measures are excluded from FEPSDE.

\textbf{Implication:} Operational rather than purely philosophical taxonomy.
\end{tcolorbox}

% DIM-8 through DIM-12 (abbreviated for space)
\begin{tcolorbox}[colback=gray!5!white, colframe=gray!75!black,
    title=Axioms DIM-8 to DIM-12: Structural Properties]

\textbf{DIM-8 (Completeness):} FEPSDE is exhaustive --- any utility source
maps to at least one dimension.

\textbf{DIM-9 (Exclusivity at Top Level):} Primary classification is mutually
exclusive (though sub-dimensions may overlap).

\textbf{DIM-10 (Temporal Stability):} Dimension structure is stable;
dimension \textit{values} change over time, structure does not.

\textbf{DIM-11 (Level Generality):} Same dimensions apply at all welfare levels
(individual through societal).

\textbf{DIM-12 (Intervention Targeting):} Each dimension can be targeted by
specific intervention types.
\end{tcolorbox}


%%%%%%%%%%%%%%%%%%%%%%%%%%%%%%%%%%%%%%%%%%%%%%%%%%%%%%%%%%%%%%%%%%%%%%%%%%%%%%%
\section{Measurement Framework}
\label{sec:what-measurement}
%%%%%%%%%%%%%%%%%%%%%%%%%%%%%%%%%%%%%%%%%%%%%%%%%%%%%%%%%%%%%%%%%%%%%%%%%%%%%%%

\begin{table}[htbp]
\centering
\caption{FEPSDE Measurement Matrix}
\label{tab:what-measurement}
\renewcommand{\arraystretch}{1.3}
\begin{tabular}{@{}clll@{}}
\toprule
\textbf{Dim} & \textbf{Objective Measure} & \textbf{Subjective Measure} & \textbf{Primary Scale} \\
\midrule
F & Income, assets, consumption & Financial satisfaction & CHF equivalent \\
E & Cortisol, heart rate variability & ESM, DRM & Affect balance \\
P & Behavioral indicators & BPNSS, mental health screens & Need satisfaction \\
S & Network size, centrality & Social support scales & Connectedness \\
D & Skill assessments, credentials & Growth orientation & Personal growth \\
X & Values-behavior alignment & MLQ, Purpose in Life & Meaning score \\
\bottomrule
\end{tabular}
\end{table}


%%%%%%%%%%%%%%%%%%%%%%%%%%%%%%%%%%%%%%%%%%%%%%%%%%%%%%%%%%%%%%%%%%%%%%%%%%%%%%%
\section{Integration with Other CORE Appendices}
\label{sec:what-integration}
%%%%%%%%%%%%%%%%%%%%%%%%%%%%%%%%%%%%%%%%%%%%%%%%%%%%%%%%%%%%%%%%%%%%%%%%%%%%%%%

\begin{tcolorbox}[colback=yellow!5!white, colframe=yellow!75!black,
    title=Integration: CORE-WHAT $\leftrightarrow$ CORE-WHO]
\textbf{What WHAT provides to WHO:}
\begin{itemize}[nosep]
\item Content of utility at each level: what dimensions exist at $L$
\item Dimension comparability across levels
\end{itemize}

\textbf{What WHAT requires from WHO:}
\begin{itemize}[nosep]
\item Bearer specification: whose utility dimensions?
\item Level-specific dimension relevance patterns
\end{itemize}
\end{tcolorbox}

\begin{tcolorbox}[colback=yellow!5!white, colframe=yellow!75!black,
    title=Integration: CORE-WHAT $\leftrightarrow$ CORE-HOW]
\textbf{What WHAT provides to HOW:}
\begin{itemize}[nosep]
\item Dimension labels for γ-matrix: $\gamma_{FE}, \gamma_{FP}, \ldots$
\item Dimension content for interaction interpretation
\end{itemize}

\textbf{What WHAT requires from HOW:}
\begin{itemize}[nosep]
\item Cross-dimension interaction values $\gamma_{dd'}$
\item Synergy/crowding specification
\end{itemize}
\end{tcolorbox}

\begin{tcolorbox}[colback=yellow!5!white, colframe=yellow!75!black,
    title=Integration: CORE-WHAT $\leftrightarrow$ CORE-WHEN]
\textbf{What WHAT provides to WHEN:}
\begin{itemize}[nosep]
\item Dimensions whose weights context modulates: $\alpha_d(\Psi)$
\item Baseline dimension structure
\end{itemize}

\textbf{What WHAT requires from WHEN:}
\begin{itemize}[nosep]
\item Context-dependent weight functions $\alpha_d(\Psi)$
\item When dimension importance shifts
\end{itemize}
\end{tcolorbox}

\begin{tcolorbox}[colback=green!5!white, colframe=green!75!black,
    title=Application Domain: Complex Problem Solving (CP)]
\textbf{CPS leverages the D(evelopmental) dimension as Competence Motivation:}
\begin{itemize}[nosep]
\item $U_D$ captures the drive to master challenges (``Kompetenzerleben'')
\item In CPS: Problem-solving success provides developmental utility
\item Intrinsic motivation for challenging problems comes from $U_D$
\end{itemize}

\textbf{Connection to Self-Determination Theory (Deci \& Ryan):}
$U_D$ corresponds to competence need satisfaction in SDT.

\textbf{Cross-Reference:} See Appendix UNMAPPED_CP (DOMAIN-CPS) for the complete 10C-CPS mapping.
\end{tcolorbox}


%%%%%%%%%%%%%%%%%%%%%%%%%%%%%%%%%%%%%%%%%%%%%%%%%%%%%%%%%%%%%%%%%%%%%%%%%%%%%%%
\section{Worked Example}
\label{sec:what-example}
%%%%%%%%%%%%%%%%%%%%%%%%%%%%%%%%%%%%%%%%%%%%%%%%%%%%%%%%%%%%%%%%%%%%%%%%%%%%%%%

\subsection{Example: Job Satisfaction Analysis}

\paragraph{Setup.}
An employee reports:
\begin{itemize}
\item High salary (F = 0.8)
\item Long commute, no lunch break (E = 0.3)
\item Micromanaged, no autonomy (P = 0.2)
\item Good team relationships (S = 0.7)
\item No learning opportunities (D = 0.2)
\item Feels work is meaningful (X = 0.8)
\end{itemize}

\paragraph{Application.}
Using equal weights and ignoring interactions:
\begin{align}
U_{\text{simple}} &= \frac{1}{6}(0.8 + 0.3 + 0.2 + 0.7 + 0.2 + 0.8) = 0.50
\end{align}

With empirical weights ($\alpha_P = 0.25$ for knowledge workers):
\begin{align}
U_{\text{weighted}} &= 0.15(0.8) + 0.10(0.3) + 0.25(0.2) + 0.20(0.7) + 0.15(0.2) + 0.15(0.8) \\
&= 0.12 + 0.03 + 0.05 + 0.14 + 0.03 + 0.12 = 0.49
\end{align}

\paragraph{Result.}
The low autonomy (P = 0.2) is particularly problematic given high $\alpha_P$
for knowledge workers. Financial compensation cannot fully offset psychological
deprivation.

\paragraph{Recommendation.}
Intervention should target P (autonomy enhancement) rather than F (more pay).


%%%%%%%%%%%%%%%%%%%%%%%%%%%%%%%%%%%%%%%%%%%%%%%%%%%%%%%%%%%%%%%%%%%%%%%%%%%%%%%
\section{Implications for Practice}
\label{sec:what-practice}
%%%%%%%%%%%%%%%%%%%%%%%%%%%%%%%%%%%%%%%%%%%%%%%%%%%%%%%%%%%%%%%%%%%%%%%%%%%%%%%

\begin{tcolorbox}[colback=green!5!white, colframe=green!75!black,
    title=Practical Implications]
\begin{enumerate}
\item \textbf{Diagnose All Six:} Assessment tools should measure all FEPSDE
      dimensions, not just satisfaction or income
\item \textbf{Target Deficits:} Interventions should address the lowest-scoring
      dimensions first (diminishing returns principle)
\item \textbf{Mind the Weights:} Different populations weight dimensions differently;
      segment-specific $\alpha_d$ values are essential
\item \textbf{Avoid Substitution Fallacy:} Money (F) cannot substitute for
      meaning (X) or autonomy (P) beyond a threshold
\end{enumerate}
\end{tcolorbox}


%%%%%%%%%%%%%%%%%%%%%%%%%%%%%%%%%%%%%%%%%%%%%%%%%%%%%%%%%%%%%%%%%%%%%%%%%%%%%%%
\section{Summary}
\label{sec:what-summary}
%%%%%%%%%%%%%%%%%%%%%%%%%%%%%%%%%%%%%%%%%%%%%%%%%%%%%%%%%%%%%%%%%%%%%%%%%%%%%%%

\begin{tcolorbox}[colback=gray!5!white, colframe=gray!75!black,
    title=Appendix UNMAPPED_C Summary: CORE-WHAT]

\textbf{Core Contribution:}
Establishes that utility comprises six irreducible dimensions (FEPSDE):
Financial, Experiential, Psychological, Social, Developmental, Existential.

\textbf{Key Formula:}
\begin{equation}
U = \sum_{d \in \text{FEPSDE}} \alpha_d \cdot u_d + \sum_{d < d'} \gamma_{dd'} \cdot u_d \cdot u_{d'}
\end{equation}

\textbf{Key Insights:}
\begin{itemize}[nosep]
\item Utility is multi-dimensional; money is one of six dimensions
\item Dimensions are irreducible; no dimension substitutes for another
\item Dimension weights vary by individual, level, and context
\item Measurement requires both objective and subjective indicators
\end{itemize}

\textbf{Next Steps:}
See Appendix UNMAPPED_B (CORE-HOW) for cross-dimension interactions ($\gamma_{dd'}$).
See Appendix UNMAPPED_V (CORE-WHEN) for context-dependent weights ($\alpha_d(\Psi)$).

\end{tcolorbox}


%%%%%%%%%%%%%%%%%%%%%%%%%%%%%%%%%%%%%%%%%%%%%%%%%%%%%%%%%%%%%%%%%%%%%%%%%%%%%%%
\section{Glossary of Symbols}
\label{sec:what-glossary}
%%%%%%%%%%%%%%%%%%%%%%%%%%%%%%%%%%%%%%%%%%%%%%%%%%%%%%%%%%%%%%%%%%%%%%%%%%%%%%%

\begin{table}[htbp]
\centering
\caption{Symbols Introduced in Appendix UNMAPPED_C}
\label{tab:what-symbols}
\renewcommand{\arraystretch}{1.3}
\begin{tabular}{@{}clp{5.5cm}c@{}}
\toprule
\textbf{Symbol} & \textbf{Name} & \textbf{Definition} & \textbf{Status} \\
\midrule
$u_F$ & Financial utility & Welfare from material resources & NEW \\
$u_E$ & Experiential utility & Welfare from hedonic experience & NEW \\
$u_P$ & Psychological utility & Welfare from need satisfaction & NEW \\
$u_S$ & Social utility & Welfare from relationships & NEW \\
$u_D$ & Developmental utility & Welfare from growth & NEW \\
$u_X$ & Existential utility & Welfare from meaning & NEW \\
$\alpha_d$ & Dimension weight & Importance of dimension $d$ & NEW \\
FEPSDE & Dimension set & $\{F, E, P, S, D, X\}$ & NEW \\
\bottomrule
\end{tabular}
\end{table}


%%%%%%%%%%%%%%%%%%%%%%%%%%%%%%%%%%%%%%%%%%%%%%%%%%%%%%%%%%%%%%%%%%%%%%%%%%%%%%%
\section{Critical Foundations and Objections}
\label{sec:what-critical}
%%%%%%%%%%%%%%%%%%%%%%%%%%%%%%%%%%%%%%%%%%%%%%%%%%%%%%%%%%%%%%%%%%%%%%%%%%%%%%%

\subsection{Objection 1: Why Six? Why Not Five or Seven?}

\begin{tcolorbox}[colback=red!5!white, colframe=red!50!black,
    title=Objection: Arbitrary Dimensionality]
\textit{``The six dimensions seem arbitrary. Why not Maslow's five? Or SDT's three?
Or a single quality-of-life index?''}
\end{tcolorbox}

\textbf{Response:}
FEPSDE emerges from synthesis of major well-being theories:
\begin{itemize}
\item Maslow (physiological/safety → F; love/esteem → S; self-actualization → D, X)
\item SDT (autonomy/competence/relatedness → P)
\item Kahneman (experienced utility → E)
\item Frankl (meaning → X)
\end{itemize}
Six is the minimum to capture all major theoretical traditions without redundancy.

\subsection{Objection 2: Cultural Specificity}

\begin{tcolorbox}[colback=red!5!white, colframe=red!50!black,
    title=Objection: WEIRD Bias]
\textit{``This taxonomy reflects Western, Educated, Industrialized, Rich, Democratic
values. Collectivist cultures may not recognize P (autonomy) or X (individual meaning).''}
\end{tcolorbox}

\textbf{Response:}
The dimensions are universal; the \textit{weights} are culturally variable.
Cross-cultural research (Diener, Suh) confirms basic psychological needs exist
universally, though expressed differently. DIM-6 (Weight Heterogeneity) accommodates
cultural variation without changing the dimensional structure.


%%%%%%%%%%%%%%%%%%%%%%%%%%%%%%%%%%%%%%%%%%%%%%%%%%%%%%%%%%%%%%%%%%%%%%%%%%%%%%%
\section{Open Issues and Research Agenda}
\label{sec:what-open}
%%%%%%%%%%%%%%%%%%%%%%%%%%%%%%%%%%%%%%%%%%%%%%%%%%%%%%%%%%%%%%%%%%%%%%%%%%%%%%%

\begin{table}[htbp]
\centering
\caption{Research Roadmap for CORE-WHAT}
\label{tab:what-roadmap}
\renewcommand{\arraystretch}{1.3}
\begin{tabular}{@{}clccl@{}}
\toprule
\textbf{\#} & \textbf{Issue} & \textbf{Priority} & \textbf{Status} & \textbf{Requirements} \\
\midrule
1 & Cross-cultural weight calibration & HIGH & Partial & Multi-country surveys \\
2 & Sub-dimension specification & HIGH & Open & Factor analysis \\
3 & Dimension measurement harmonization & MEDIUM & Ongoing & Psychometric work \\
4 & Temporal dynamics of dimensions & MEDIUM & Open & Panel data \\
5 & Neural correlates per dimension & LOW & Open & Neuroscience \\
\bottomrule
\end{tabular}
\end{table}


%%%%%%%%%%%%%%%%%%%%%%%%%%%%%%%%%%%%%%%%%%%%%%%%%%%%%%%%%%%%%%%%%%%%%%%%%%%%%%%
\section{References}
\label{sec:what-references}
%%%%%%%%%%%%%%%%%%%%%%%%%%%%%%%%%%%%%%%%%%%%%%%%%%%%%%%%%%%%%%%%%%%%%%%%%%%%%%%

\begin{tcolorbox}[colback=blue!5!white, colframe=blue!75!black]
\textbf{Master Bibliography:} See \texttt{bcm\_master.bib}
\end{tcolorbox}

\subsection{Key References}

\begin{itemize}[nosep]
\item \textbf{Experienced Utility:} \citet{kahneman1999wellbeing}
\item \textbf{Self-Determination:} \citet{ryan2000self}
\item \textbf{Meaning:} \citet{frankl1959search}, \citet{steger2006meaning}
\item \textbf{Psychological Well-Being:} \citet{ryff1989happiness}
\item \textbf{Social Capital:} \citet{putnam2000bowling}
\end{itemize}

\nocite{bcm_master}

%%%%%%%%%%%%%%%%%%%%%%%%%%%%%%%%%%%%%%%%%%%%%%%%%%%%%%%%%%%%%%%%%%%%%%%%%%%%%%%
% FOOTER
%%%%%%%%%%%%%%%%%%%%%%%%%%%%%%%%%%%%%%%%%%%%%%%%%%%%%%%%%%%%%%%%%%%%%%%%%%%%%%%

\vspace{1em}
\hrule
\vspace{0.5em}

\noindent\textit{Cross-references:}
Appendix UNMAPPED_GLS (Glossary),
Appendix UNMAPPED_AAA (CORE-WHO),
Appendix UNMAPPED_B (CORE-HOW),
Appendix UNMAPPED_V (CORE-WHEN),
Appendix UNMAPPED_BBB (CORE-WHERE),
Appendix UNMAPPED_AU (CORE-AWARE).

\noindent\textit{Master Bibliography:} \texttt{bcm\_master.bib}

% =============================================================================
% END OF APPENDIX C
% =============================================================================
